\documentclass[12pt]{article}
\usepackage[utf8]{inputenc}
\usepackage[T1]{fontenc}
\usepackage{hyperref}
\title{Ubang Language Resources for K-12 Curriculum Development}
\author{Compiled by Claude Sonnet 5 (claude-sonnet-5) for Andr\'as Kornai}
\date{}
\begin{document}
\maketitle
\section{Introduction}
Ubang (ISO 639-3 uba) is a Bendi language of Obudu LGA, Cross River State, Nigeria, with about 11,000 speakers (2013). It is known for gender-based diglossia, in which men and women use different word forms for the same things. A documentation project at the University of Ibadan, funded by ELDP, classes it as endangered, citing emigration, contact with Nigerian Pidgin and English, and weaker transmission among younger speakers. Joshua Project lists the villages as Ofambe, Okiro and Okweriseng.
\section{Primary Educational Resources}
No primers, textbooks or school programmes were found.
\section{Linguistic and Reference Materials}
\begin{itemize}
\item Wikipedia, Ubang language: \url{https://en.wikipedia.org/wiki/Ubang_language}
\item ELDP archive, Documentation of Ubang Gender Diglossia: \url{https://eldp.access.preservica.com/dk0492/}
\item Glottolog reference, Ugwu, Ifeanyi Aaron 2018: \url{https://glottolog.org/resource/reference/id/582092}. I did not verify its content.
\end{itemize}
\section{Religious and Community Resources}
\begin{itemize}
\item Joshua Project, Ubang in Nigeria: \url{https://joshuaproject.net/people_groups/15678}. It notes that a full Bible translation is still needed.
\end{itemize}
\section{Digital and Audio Resources}
The ELDP documentation collection above is the only candidate. I did not inspect its contents, so I do not know whether it includes audio.
\section{Curriculum Development Context}
The documentation project is the natural foundation for any curriculum. Materials would need to deal with the male and female speech forms, which the sources report as mutually intelligible. Teaching resources would have to be created from scratch.
\end{document}